\begin{proofbox}
\textbf{\textcolor{CProof}{思路提示}}

把 $A_n^m$ 和 $A_n^{m-1}$ 都按排列数定义展开，比较两式相差的最后一个因子即可。

\medskip
\textbf{\textcolor{CProof}{正式推导}}
\begin{description}[style=nextline, leftmargin=2.8em, labelsep=0.5em, itemsep=0.55em]
    \item[\textbf{步骤一（定义展开）：}] 写出 $A_n^m$ 的定义。
    \[
    A_n^m=n(n-1)(n-2)\cdots(n-m+1).
    \]
    \par\textbf{理由：} 从 $n$ 个不同元素中取 $m$ 个排列，依次有 $n,n-1,\ldots,n-m+1$ 种选择。

    \item[\textbf{步骤二（写出 $A_n^{m-1}$）：}] 写出 $A_n^{m-1}$ 的定义。
    \[
    A_n^{m-1}=n(n-1)(n-2)\cdots(n-m+2).
    \]
    \par\textbf{理由：} 从 $n$ 个不同元素中取 $m-1$ 个排列，最后一个因子到 $n-m+2$ 为止。

    \item[\textbf{步骤三（比较乘积）：}] $A_n^m$ 比 $A_n^{m-1}$ 多出最后一个因子 $n-m+1$。
    \[
    \begin{aligned}
    A_n^m
    &=n(n-1)\cdots(n-m+2)(n-m+1)\\
    &=(n-m+1)\bigl[n(n-1)\cdots(n-m+2)\bigr]\\
    &=(n-m+1)A_n^{m-1}.
    \end{aligned}
    \]
    \par\textbf{理由：} 这里使用的是乘法交换律与结合律，将多出的因子单独提出。
\end{description}

\medskip
\textbf{\textcolor{CProof}{结论回扣}}

因此，当 $n,m$ 为正整数且 $2\le m\le n$ 时，
\[
A_n^m=(n-m+1)A_n^{m-1}.
\]
证明完毕。
\end{proofbox}
